%% LyX 2.5.1 created this file.  For more info, see https://www.lyx.org/.
%% Do not edit unless you really know what you are doing.
\documentclass[letterpaper,english,reprint, aps]{revtex4-1}
\usepackage[T1]{fontenc}
\usepackage[latin9]{inputenc}
\setcounter{secnumdepth}{3}
\usepackage{amsmath}
\usepackage{amssymb}

\makeatletter

%%%%%%%%%%%%%%%%%%%%%%%%%%%%%% LyX specific LaTeX commands.
\pdfpageheight\paperheight
\pdfpagewidth\paperwidth


\makeatother

\usepackage{babel}
\begin{document}
\title{Auxiliary Field Quantum Monte Carlo Simulations of Non-hermitian Bosonic
Systems Without the Sign Problem}
\author{Yin-Kai Yu}
\affiliation{University of Chinese Academy of Sciences, Beijing 100049, China}
\author{Zhong-Chao Wei}
\email{zcwei@thp.uni-koeln.de}

\affiliation{Qingdao Binhai Univeristy, Qingdao 266555, China}
\author{Lei Wang}
\email{wanglei@iphy.ac.cn}

\affiliation{Beijing National Laboratory for Condensed Matter Physics and Institute
of Physics,~\\
Chinese Academy of Sciences, Beijing 100190, China}
\affiliation{Songshan Lake Materials Laboratory, Dongguan, Guangdong 523808, China}
\author{Zi-Xiang Li}
\email{zixiangli@iphy.ac.cn}

\affiliation{Beijing National Laboratory for Condensed Matter Physics and Institute
of Physics,~\\
Chinese Academy of Sciences, Beijing 100190, China}
\affiliation{University of Chinese Academy of Sciences, Beijing 100049, China}
\date{\today}
\begin{abstract}
This is a note on sign-problem-free non-Hermitian models in bosonic
auxiliary field quantum Monte Carlo.
\end{abstract}
\keywords{sign problem, non-hermitian system, quantum Monte Carlo}
\maketitle

\section{Introduction}

The concept of contraction semigroup was used to study the sign problem
of quantum lattice fermion models\citep{wei_semigroup_2024}. Now
we extend this method to bosonic auxiliary field quantum Monte Carlo
(BAFQMC). Instead of hermitian systems, this method is more closely
related to non-hermitian bosonic Hamiltonians.

\section{The trace formula in AFQMC}

In BAFQMC algorithms, interaction terms are decoupled into bosonic
quadratic forms by introducing auxiliary fields. The action in terms
of auxiliary fields can be treated with random sampling numerically.
For the bosonic case, the sampling weight can be expressed by\citep{thomas_character_2008}
\begin{align}
p & =\mathrm{tr}\left(e^{h_{1}}e^{h_{2}}\cdots e^{h_{k}}\right)\nonumber \\
 & =\det\left(I-e^{A_{1}}e^{A_{2}}\cdots e^{A_{k}}\right)^{-1/2}.\label{eq:TraceFormulaBoson}
\end{align}
Here $h_{i}=-\frac{1}{2}\chi^{T}\Omega A_{i}\chi$ ($i=1,2,\dots k$)
are a set of bosonic quadratic forms. We have used the basis $\chi=\left(\begin{array}{c}
a_{1}\\
a_{2}\\
\vdots\\
a_{N}\\
a^{+}_{1}\\
a^{+}_{2}\\
\vdots\\
a^{+}_{N}
\end{array}\right)$, where the bosonic creation and annihilation operators satisfy the
canonical commutation relations $\left[a_{i},\ a_{j}\right]=0$, $\left[a^{+}_{i},\ a^{+}_{j}\right]=0$,
$\left[a_{i},\ a^{+}_{j}\right]=\delta_{ij}$($i,j=1,2,\dots N$),
and $\Omega=\left(\begin{array}{cc}
0 & I_{N}\\
-I_{N} & 0
\end{array}\right)$. In principle the coefficient matrices $A_{i}$ could be arbitrary
complex square matrices $A$ that satisfies $\Omega A+A^{T}\Omega=0$.

If we do not put any restriction on $h_{i}$, the sampling weight
$p$ could be non-positive, for one thing the coefficients are complex,
and for another the spin or the metapletic representation is two-valued.
Statistical sampling methods usually cannot obtain desired physical
quantities with useful accuracy at reasonable cost under those circumstances.
This is the notorious sign problem in AFQMC methods.

We want to emphasize that the second equal sign in Eq.~(\ref{eq:TraceFormulaBoson})
should be interpreted as the regularized sampling weight instead of
the bare sampling weight{[}cite our work{]}.

\section{Semigroups and sign-problem-free AFQMC simulations }

For the semigroup defined by
\begin{align}
J^{T}_{1}AJ_{1} & =\bar{A},\\
J_{2}\Omega A+A^{\dagger}J_{2}\Omega & \leq0,
\end{align}
where (i)$J^{T}_{1}\Omega J_{1}=-\Omega$, $J^{2}_{1}=I_{2N}$, $J^{T}_{2}\Omega J_{2}=-\Omega$,
$J^{2}_{2}=I_{2N}$, $\left\{ J_{1},\ J_{2}\right\} =0$; or (ii)$J^{T}_{1}\Omega J_{1}=-\Omega$,
$J^{2}_{1}=-I_{2N}$, $J^{T}_{2}\Omega J_{2}=-\Omega$, $J^{2}_{2}=I_{2N}$,$\left\{ J_{1},\ J_{2}\right\} =0$;
we have $p>0$.

The proof follows the scheme in Ref.~\citep{wei_semigroup_2024}.
Firstly we restrict the range of parameters by a $Z_{2}$ anti-linear
symmetry represented by $J_{1}$ to make the sampling weight $p$
real. Then we observe that $p$ never goes to $\pm\infty$ for strict
contraction elements inside some contraction semigroups, where the
metric of the contraction semigroups is given by $J_{2}\Omega$, $J_{2}$
represents another $Z_{2}$ anti-linear symmetry. The non-strict contraction
elements can be viewed as some limits of strict contraction elements.
These two conditions together ensure that $p^{-1}$ is non-negative
as a continuous function of the coefficients. The consistency of the
two conditions requires $\left\{ J_{1},J_{2}\right\} =0$.

\section{Sign-problem-free non-hermitian models}

For the bosonic case, we study the non-hermitian models with anti-hermitian
single-body terms $H_{K}=-H^{\dagger}_{K}$ and hermitian interaction
terms $H_{U}=H^{\dagger}_{U}$.

For case (i), we could choose $J_{1}=\left(\begin{array}{cc}
I_{N} & 0\\
0 & -I_{N}
\end{array}\right)$, $J_{2}=\left(\begin{array}{cc}
0 & I_{N}\\
I_{N} & 0
\end{array}\right)$. We have
\begin{equation}
H_{K}=a^{+T}Ba,
\end{equation}
\begin{equation}
H_{U}=U\sum_{\left\langle i,j\right\rangle }\left(a^{+}_{i}a_{j}-a^{+}_{j}a_{i}\right)^{2},
\end{equation}
where $B=\bar{B}=-B^{T}$, $U\leq0$.

For case (ii), we could choose $J_{1}=\left(\begin{array}{cccc}
0 & I_{N/2} & 0 & 0\\
-I_{N/2} & 0 & 0 & 0\\
0 & 0 & 0 & -I_{N/2}\\
0 & 0 & I_{N/2} & 0
\end{array}\right)$, $J_{2}=\left(\begin{array}{cc}
0 & I_{N}\\
I_{N} & 0
\end{array}\right)$. Let $b_{i}$, $b^{+}_{i}$($i=1,2,\dots N/2$) and $c_{i}$, $c^{+}_{i}$($i=1,2,\dots N/2$)
be two sets of bosonic annihilation/creation operators. We have

\begin{equation}
H_{K}=b^{+T}B_{11}b+c^{+T}\bar{B}_{11}c+b^{+T}B_{12}c-c^{+T}\bar{B}_{12}b,
\end{equation}
\begin{equation}
H_{U}=U\sum_{i}\left(b^{+}_{i}b_{i}-c^{+}_{i}c_{i}\right)^{2},
\end{equation}
where $B_{11}=-B^{\dagger}_{11}$, $B_{12}=B^{T}_{12}$, $U\geq0$.
When $B_{12}\neq0$, this model does not have reflection positivity,
but it is still sign-problem-free.
\begin{acknowledgments}
We would like to thank...
\end{acknowledgments}

\bibliographystyle{apsrev4-1}
\bibliography{SGNH}
\newpage{}

\appendix

\section*{Supplemental Material}

\section{appendix}
\end{document}
